The first fact follows from \(\text{thm:adaptive-rootn-minimax}\).  That theorem selects and fixes one once-for-all bank supplied by \(\text{lem:finite-contour-bank}\) and constants independent of \(n\), and on every stated index proves
\[
 c/n\le\mathfrak R_n^{\mathrm{NG}}\le C/n.
\]
It proves the upper bound using the measurable statistic associated with that selected bank and does not construct the bank effectively.  Inspection of \(\text{def:non-gaussian-class}\) shows that membership uses \(\bar g_n\) and \(\varepsilon_{1,n}\), but no supplied outcome code or outcome-code radius.  Inspection of \(\text{def:adaptive-contour-estimator}\) likewise shows that, conditional on the selected bank, the statistic uses \(Y\) and \(T-\bar g_n(X)\), but not \(\widehat q_n,\bar q_n\), or \(\varepsilon_{2,n}\).  This proves the fixed-separated non-Gaussian assertion, including fixed nonzero treatment-code error.

For the separate Gaussian class, \(\text{prop:bounded-outcome-gaussian-degeneracy}\) proves that every member has \(\theta_0(P)=0\).  When that class is nonempty, the same proposition and \(\text{def:minimax-risks}\) give
\[
 \mathfrak R_n^{\mathrm G}=\mathfrak M_{n,1-\gamma}^{\mathrm G}=0.
\]
Its argument uses the simultaneous exact conditional-mean equation, bounded observed outcome, and nondegenerate Gaussian treatment residual: a bounded outcome has a bounded conditional mean, whereas a nonzero coefficient multiplying the Gaussian residual has unbounded support.  Hence the incompatibility is internal to those simultaneous source assumptions, and no positive Gaussian lower bound is used.  Since the conclusions concern separately defined classes and neither inclusion nor equality has been established, they are separate facts and imply neither a Gaussian-versus-non-Gaussian minimax dichotomy nor a shrinking-separation frontier.